% !PRINTED 221
\weberpage{221}
%<<CONT
first invariant of the form biquadratic with respect to the variables $\xi_1, \xi_2$, namely (Bd.\ I, \S.\,64):

\weq{3}{u_2^2 - 3\,u_1u_3 + 12\,u_0u_4.}

But since a form of degree 16 cannot be represented as a product of forms of degrees 12, 20 and 30, this covariant must vanish identically. Now here

\[
u_0 = 11\,.\,10\,.\,9\,.\,8\,x_1^7x_2 + 6\,.\,5\,.\,4\,.\,3\,m\,x_1^2x_2^6,
\]
\[
u_1 = 4\,.\,11\,.\,10\,.\,9\,x_1^8 + 4\,.\,6\,.\,5\,.\,4\,.\,6\,m\,x_1^3x_2^5,
\]
\[
u_2 = 6^3\,.\,5^2\,m\,x_1^4x_2^4,
\]
\[
u_3 = -4\,.\,11\,.\,10\,.\,9\,x_2^8 + 4\,.\,6\,.\,5\,.\,4\,.\,6\,m\,x_2^3x_1^5,
\]
\[
u_4 = -11\,.\,10\,.\,9\,.\,8\,x_2^7x_1 + 6\,.\,5\,.\,4\,.\,3\,m\,x_2^2x_1^6,
\]

and if we form from these the covariant (3) and set the coefficient of $x_1^8x_2^8$ equal to 0, we obtain $m = \pm 11$. Both numbers are here admissible. One choice is transformed into the other by replacing $x_1$ with $-x_1$, and hence by a transformation of the group. We therefore set:

\weq{4}{f(x_1,x_2) = x_1x_2(x_1^{10} + 11\,x_1^5x_2^5 - x_2^{10}).}

The ten five-fold poles still missing are therefore obtained by solution of the equation

\[
x^{10} + 11\,x^5 - 1 = 0,
\]

from which one finds the following values:

\[
x^5 = \frac{-11 \pm 5\sqrt5}{2} = \left(\frac{-1 \pm \sqrt5}{2}\right)^5
\]

If we therefore set

\weq{5}{\begin{aligned}
\omega &= \varepsilon + \varepsilon^4 = 2\cos\frac{2\pi}{5} = \frac{-1+\sqrt5}{2} \\
\omega' &= \varepsilon^2 + \varepsilon^3 = 2\cos\frac{4\pi}{5} = \frac{-1-\sqrt5}{2},
\end{aligned}}

so that

\weq{6}{\omega^2 + \omega = 1, \quad \omega + \omega' = -1, \quad \omega\omega' = -1}

holds, then the five-fold poles other than 0 and $\infty$ are:

\weq{7}{\xi = \varepsilon^\lambda\omega, \quad \varepsilon^\lambda\omega' = \frac{-\varepsilon^\lambda}{\omega}, \quad \lambda = 0, 1, 2, 3, 4.}

Now we apply the theorem \S.\,53, 2., according to which, if $a, a'$ are two poles of a substitution $\Theta$, and $b$ is a pole conjugate to $a$, there must exist a substitution $\chi$ in the group such
%CONT>>
